% Foundations of Physics - Cover Letter
% FUND-FOP-047: Golden Ratio φ, e, π and Fine Structure Constant α: Collapse Breathing Proportions

\documentclass[12pt]{letter}
\usepackage{geometry}
\usepackage{hyperref}
\usepackage{graphicx}
\usepackage{fancyhdr}

\geometry{a4paper, margin=1in}

\begin{document}

\begin{letter}{Editor-in-Chief\\
Foundations of Physics\\
Springer Nature}

\opening{Dear Editor-in-Chief,}

We are pleased to submit our manuscript entitled ``Golden Ratio $\phi$, $e$, $\pi$ and Fine Structure Constant $\alpha$: Collapse Breathing Proportions'' for consideration for publication in Foundations of Physics.

This paper presents novel research on the mathematical relationships between four fundamental constants: the Golden Ratio ($\phi$), Euler's number ($e$), Pi ($\pi$), and the Fine Structure Constant ($\alpha$). Using the Universe Ontology framework, we demonstrate that these seemingly disparate constants form a unified system of proportions that we term ``collapse breathing proportions,'' which govern fundamental reality structures through FLIP-XOR-SHIFT operations.

The significance of our work includes:

\begin{enumerate}
\item Derivation of exact mathematical relationships between $\phi$, $e$, $\pi$, and $\alpha$ using information-theoretic principles
\item Introduction of ``collapse breathing proportions'' as a unifying concept for understanding constant relationships
\item Development of a theoretical explanation for the specific value of the fine structure constant
\item Numerical verification of these relationships to high precision
\item Exploration of the philosophical and physical implications of these findings
\end{enumerate}

Our research contributes to the ongoing discourse about the nature of physical constants and offers a new perspective on their origin as manifestations of underlying information-theoretic principles rather than arbitrary values. We believe this work aligns well with Foundations of Physics' tradition of publishing foundational research that challenges conventional understanding and proposes innovative theoretical frameworks.

All authors have approved the manuscript and agree with its submission to Foundations of Physics. This manuscript has not been published elsewhere and is not under consideration by any other journal.

We suggest the following individuals who have expertise in this area and could serve as potential reviewers:

\begin{enumerate}
\item Prof. Carlo Rovelli, Aix-Marseille University, carlo.rovelli@univ-amu.fr - Expert in quantum gravity and foundations of physics
\item Dr. Gerald Gabrielse, Northwestern University, gabrielse@northwestern.edu - Leading expert on precision measurements of fundamental constants
\item Prof. Max Tegmark, Massachusetts Institute of Technology, tegmark@mit.edu - Expert in mathematical physics and information theory approaches
\item Prof. David H. Bailey, University of California, Davis, dhbailey@lbl.gov - Authority on high-precision computation and mathematical constants
\end{enumerate}

We would prefer to avoid researchers directly affiliated with competing information-theoretic frameworks as reviewers due to potential conflicts of interest.

Thank you for your consideration of our manuscript. We look forward to your response.

\closing{Sincerely,}

\vspace{0.5cm}
Haobo Ma (Corresponding Author)\\
AELF PTE LTD.\\
8 Marina Blvd, \#14-02, Marina Bay Financial Centre Tower 1\\
Singapore 018981\\
Email: auric@aelf.io\\
ORCID: 0009-0008-4944-977X

\vspace{0.5cm}
Wen Niu\\
AELF PTE LTD.\\
Email: ada@aelf.io\\
ORCID: 0009-0006-3349-0298

\end{letter}

\end{document} 